% Replacement for the "Dyna iteration" paragraph of \subsection{OGBench Cube}
% (app:results-ogbench), 2026-09-25. The previous version described Dyna v1: collection
% by the PRE planner at h25 only, an unanchored 2-epoch fine-tune, and the finetuned
% model "reused as-is for h100". The unified campaign changed the collector, the
% objective, the number of rounds and what is reported; the h100 claim is no longer made.

\paragraph{Setup.}
Every cell is $50$ episodes per evaluation seed. Success follows the benchmark's
cube-placement criterion; the no-op floors ($56.0$ at $h25$, $45.3$ at $h100$) are
measured by executing zero actions under the identical protocol.

\paragraph{Dyna iterations.}
Each round collects on-policy episodes with the \emph{deployed} planners of the previous
round -- all six seeds of the cell's locked (teacher, actor) pair, the same actors we
report -- on training-split tasks (episodes $0$--$7999$, no termination at goal). We
collect at \emph{both} evaluation horizons, $1{,}800$ episodes at $h25$ and $1{,}800$ at
$h100$, and mix them $50{:}50$ with the original expert data, outcome-labeled. The world
model is then fine-tuned for one epoch at LR $10^{-5}$ from the previous round's model,
with a latent anchor $\lVert \mathrm{enc}_{\mathrm{ft}}(x) - \mathrm{enc}_{\mathrm{prev}}(x)\rVert^2$
(weight $1.0$) added to the objective on every frame of the mixture; the encoder is
trained, not frozen. The anchor exists because the critic and actor are rebuilt from
scratch on the new latent space, so drift in regions the on-policy data does not cover
would silently invalidate what the base model knew. PLDM is fine-tuned under its own
native objective rather than LeWM's. Latent caches and the offline critic are rebuilt
under the fine-tuned model, and RP1 is retrained \emph{at the cell's locked
configuration} -- the (teacher budget, actor step) pair selected in
Sec.~\ref{app:results-general} -- never re-selected, so the comparison isolates the
world model. Round $k{+}1$ repeats this with round $k$'s planners.

We run two rounds and report, per cell, the round with the better validation score:
round~$1$ for LeWM ($93.00$ against round~2's $92.75$) and round~$2$ for PLDM ($90.33$
against $83.25$). We report Dyna at $h25$ only. The two horizons fail on disjoint task
sets -- $9$ of $150$ at $h25$ against $17$ of $150$ at $h100$, different segments of the
same episodes -- and across rounds the fine-tune repaired $3$--$6$ of the $h25$ core
tasks but $1$ of the $h100$ core, leaving $h100$ changes inside seed noise. Collection
horizon is not the explanation: we collect at both offsets and the $h100$ core still does
not move, consistent with long-horizon failures being errors of execution over a long
contact sequence rather than of one-step dynamics.
